Planning:


-> Talk about how remote sensing temperature of rivers with TIR images is done more and more often, as it poses advantages over other techniques.
-> higher spatial resolution possible

-> There are however also unique challenges
-> TIR sensors measure the xxx which is often not equal to the kinematic temperature as the xxx temperature is a function of the kinematic temperature, which for instance also depends on the emissivity $\epsilon$ of the material. This emissivity is often not known.

-> The TIR sensors used also introduce a source of error. There are cooled and uncooled sensors. Cooled sensors generally yield measurement results with a higher signal to noise ratio, as they can controll the sensor temperature, which is a primary source of noise. Uncooled sensors on the other hand lack this capability and are therefore subject to so called "temperature drift". As the sensor operates, the sensor parts like for instance the lense heat up. This thermal radiadtion stemming from the sensor itself are also picked up, which increases the measured temperature of the land surface. This leads to a steady increase of the measured temperature of the land surface, even tough the kinematic temperature of the surface might not increase. Further, the sensor can also be cooled down trough wind, brushing the sensor and therefore striping away heat. Then the sensor readings slowly decrease over time. This drift has to be accounted for and different hardware aswell software based correction methods exist and are included into the sensor to various degrees. An example is the NUC shutter present in a multitude of different cameras, however there are also other tools like the xy and the xxy.

Studies have shown, that even with these correction methods, significant temperature drift can still accumulate over the flight mission. It is therefore advisable to correct the temperature drift in post processing, which is often done using temperature loggers located on the land surface. Different correciton mechanisms are then available to correct for temperature drift and generally to express the relationship between the xxx temperature and the kinetic temperature.

Deploying the loggers is generally costly and time demanding. It is therefore of interest to achieve high quality correction results, with as little loggers as possible.



-> different radiometric methods for correcting the temperature drift in the data of this project will be tested. It will also be tested, which method performs best with as view loggers as well:



Planning refined:

-> Thermal heterogenity in rivers is vital for biodiversity and especially the presence of so called coldwater refugia is of great importance for the survival of fish with the respect of increasing temperatures of river systems.
-> Mapping out these coldwater refugia is thereofre of great importance. This can only be marginally achieved with the current technologies and the application of TIR remote sensing is becoming more and more common for Temperature mapping of streams.

-> As part of a large cantonal project, the ZHAW has carried out TIR collection missions for over 100 km of river transect of the emme river aswell as selected tributaries in the canton of bern. A workflow was then derived to process the optained TIR images and finally derive the location and extend of coldwater spots in the emme river aswell as longitudinal and queere temperature gradients. (swap this). As part of this workflow, a temperature correction step is applied, where Brightes Temperature of the TIR data is corrected to reflect the kinetic temperature of the Water. This correction mechanism is carried out in two steps: 


The highly overlapping TIR images are beeing stiched together using the Pix4DMapper sorftware. Every pixel in the Resulting TIR orthophoto has a range of different temperature Values associated to it due to the high overlapp of the singular TIR images. From these values, different per pixel statistics where like the tempearature mean, the range etc where calculated. The pixels located on the middle line are then extracted for every flight transect and compared to the middle line pixels of the other flight transects. This is used to make messurement distortions due to e.g. temperature drift visible. Distorded pixels or even whole Flight transects are then removed. (tonnolla et all 2021). 

Then a transect based correction method is applied, where all water pixels associated with a specific transect get corrected linearly based on the logger temperature measured in this transect. 



In this BSc Project, other algorithms, that can be used to address the temperature drift of the sensor aswell as the changing environmental conditions like cloud and wind in TIR images will be Explored. Speciffically, relative radiometric correction mechanisms, that can be applied to individual TIR images will be applied. The results of the different radiometric correction methods will then be compared amongst eachother aswell to the currently applied correction mechanism. The RMSE of the methods will be recorded and compared. Further an analysis on the amount of loggers with respect to the quality of the correction will be determined and compared between the methods. This, since deploying loggers for thermal correction is cosstly and demands a lot of effort.


